% SHARD-ID: AISM-26-HCB4-CANONICAL-GRAM
% SHARD-TITLE: The canonical corner Gram estimate
% SHARD-SUMMARY: Reproduces lem-hcb4-canonical-gram, the af-validated two-sided comparison of the ambient corner norm with the norm of the canonical identification.
% SHARD-SUMMARY: Explains the exact Gram scalarisation, the two uniform inputs it is fed, the scalar square-root passage, and the adjoint transfer to the row corner.
% SHARD-KEYWORDS: lem-hcb4-canonical-gram, af validated, canonical corner identification, Gram scalarisation, compressed unit, row transfer
\section{The canonical corner Gram estimate}\label{sec:hcb4-canonical-gram}

\begin{lemma}[\texttt{lem-hcb4-canonical-gram}]\label{lem:hcb4-canonical-gram}
There are universal $C_J<\infty$ and $\comb_J>0$ such that every H-CB datum with
$\comb\le \comb_J$, every $n\ge1$ and every $Z$ in either special $P,Q$ corner satisfy
\begin{equation}\label{eq:cg-main}
  (1-C_J\comb)\,\nrm{Z}\;\le\;\bigl\lVert J_n(Z)\bigr\rVert\;\le\;(1+C_J\comb)\,\nrm{Z}.
\end{equation}
\emph{Transcription note.} The \emph{special} corners are the two involving the
one-dimensional projection of the datum: $Z\in\Amp{n}{\Cnr{P}{Q}}$ with $J_n=J_{P,Q,n}$
the canonical column identification, or $Z\in\Amp{n}{\Cnr{Q}{P}}$ with $J_n=J_{Q,P,n}$ the
canonical row identification, both fixed by
\texttt{def-canonical-corner-identifications}. The two sides of \eqref{eq:cg-main} carry
\emph{different} norms: $\nrm{Z}$ is the ambient operator-space norm of the amplified
corner, while $\lVert J_n(Z)\rVert$ is the operator norm of the resulting map between
amplified column spaces --- the Euclidean norm of \texttt{def-ha-map}. The content of the
lemma is therefore that the canonical identification, which is pure notation and carries
no estimate of its own, is a $1\pm O(\comb)$ isometry between those two norms, uniformly
in $n$.
\end{lemma}

\noindent\textbf{Registry contract (verbatim).}
\contractquote{Canonical corner Gram estimate: there are universal C_J < infinity and e_J > 0 such that every H-CB datum with e <= e_J, every n >= 1, and every Z in either special P,Q corner satisfy (1-C_J*e)||Z|| <= ||J_n(Z)|| <= (1+C_J*e)||Z||.}

\paragraph{Proof (prose account of the af-validated tree).}
The tree proves the column case and transfers the row case to it. The column case is one
exact identity fed by two uniform inputs, followed by a scalar computation; there is no
analysis beyond the two imports.

\emph{Exact Gram scalarisation.} Fix an admissible datum, $n\ge1$ and
$Z\in\Amp{n}{\Cnr{P}{Q}}$, and put $G:=J_{P,Q,n}(Z)\dg J_{P,Q,n}(Z)\in\Mat{n}$. Reading
the defining formula of \texttt{def-canonical-corner-identifications} against the
column-Hilbert inner-product displays of \texttt{def-column-hilbert-corner} and the
compressed-product display, the compressed square of $Z$ is computed entrywise:
\begin{equation}\label{eq:cg-entry}
  \bigl[Z\dg\cpr Z\bigr]_{ij}
  \;=\;\sum_k\ipo{Z_{ki}}{Z_{kj}}\,\cunit{Q}
  \;=\;G_{ij}\,\cunit{Q},
  \qquad\text{that is}\qquad
  Z\dg\cpr Z\;=\;G\otimes\cunit{Q}.
\end{equation}
This is an identity, not an estimate. Two norm facts turn it into one. The Ruan
scalar-tensor invariance among the operator-space matrix-norm axioms gives
$\nrm{G\otimes\cunit{Q}}=\nrm{G}\,\nrm{\cunit{Q}}$, and the Hilbert-space Gram identity
gives $\nrm{G}=\lVert J_{P,Q,n}(Z)\rVert^{2}$. Hence
\begin{equation}\label{eq:cg-scalar}
  \bigl\lVert Z\dg\cpr Z\bigr\rVert
  \;=\;\bigl\lVert J_{P,Q,n}(Z)\bigr\rVert^{2}\,\nrm{\cunit{Q}} .
\end{equation}
Everything after this point is bookkeeping on two scalars: the compressed square on the
left and the compressed unit on the right.

\emph{The two uniform inputs.} There are universal $A,B<\infty$, valid below a universal
threshold, with
\begin{equation}\label{eq:cg-inputs}
  \bigl\lvert\,\lVert Z\dg\cpr Z\rVert-\nrm{Z}^{2}\,\bigr\rvert\le A\,\comb\,\nrm{Z}^{2},
  \qquad
  \bigl\lvert\,\nrm{\cunit{Q}}-1\,\bigr\rvert\le B\,\comb .
\end{equation}
For the first, the $\ecs$-$C^*$ norm axioms of \texttt{def-extended-epsilon-cstar-algebra}
place the \emph{ambient} square $\nrm{Z\dg Z}$ between $(1-\ecs)\nrm{Z}^{2}$ and
$(1+\ecs)\nrm{Z}^{2}$ at level $n$, while
\texttt{lem-compcb-rectangular-product} (\S\ref{sec:compcb-rectangular-product}) applied to
the compatible amplified pair $(Z\dg,Z)$ gives
$\lVert Z\dg\cpr Z-Z\dg Z\rVert\le C_r\comb\nrm{Z}^{2}$; since $\ecs\le\comb$, the reverse
triangle inequality yields the first bound with $A=C_r+1$, uniformly in $n$. For the
second, \texttt{lem-compcb-compressed-unit-norm}
(\S\ref{sec:compcb-compressed-unit-norm}) gives $\lvert\nrm{\cunit{Q}}-1\rvert\le
C_u\comb$ once $Q$ is known to be \emph{nonvanishing}; the datum states only that $Q$ is
one-dimensional, and the gap is closed as elsewhere in this tier by the registered bridge
$\dim\Cnrs{Q}=1\Rightarrow\Cnrs{Q}\ne0\Rightarrow Q$ nonvanishing. So $B=C_u$.

\emph{The scalar passage.} Assume $Z\ne0$ (for $Z=0$ both sides of \eqref{eq:cg-main}
vanish) and put $x:=\lVert J_{P,Q,n}(Z)\rVert/\nrm{Z}$. Dividing \eqref{eq:cg-scalar} by
$\nrm{Z}^{2}$ gives $x^{2}\nrm{\cunit{Q}}=\lVert Z\dg\cpr Z\rVert/\nrm{Z}^{2}$, so
\eqref{eq:cg-inputs} confines
\[
  x^{2}\;\in\;\Bigl[\tfrac{1-A\comb}{1+B\comb},\ \tfrac{1+A\comb}{1-B\comb}\Bigr].
\]
Take $\comb_{\mathrm{col}}$ no larger than the two external thresholds, $1/(2\max\{B,1\})$
and $1/\max\{A+B,1\}$. Then the denominators are bounded away from $0$ and the two ends
simplify to $x^{2}\ge1-(A+B)\comb$ and $x^{2}\le1+2(A+B)\comb$. Finally
$\sqrt{1-t}\ge1-t$ for $t\in[0,1]$ and $\sqrt{1+t}\le1+t/2$ for $t\ge0$ convert these into
\begin{equation}\label{eq:cg-col}
  1-(A+B)\comb\;\le\;x\;\le\;1+(A+B)\comb,
\end{equation}
which is \eqref{eq:cg-main} for the column corner with $C_{\mathrm{col}}=A+B$. The
asymmetric factor $2$ in the upper bound for $x^{2}$ is exactly what the sharper
square-root inequality absorbs, so the same constant serves both sides.

\emph{Row transfer.} For $Z\in\Amp{n}{\Cnr{Q}{P}}$ put $W:=Z\dg\in\Amp{n}{\Cnr{P}{Q}}$.
The self-adjoint operator-space axiom gives $\nrm{W}=\nrm{Z}$ at every level;
\texttt{def-canonical-corner-identifications} \emph{defines} the row map by
$J_{Q,P,n}(Z)=J_{P,Q,n}(W)\dg$; and Hilbert-space adjunction preserves the operator norm.
Applying the column case to $W$ therefore returns both row inequalities with the same
constants, so $C_J=C_{\mathrm{col}}$ and $\comb_J=\comb_{\mathrm{col}}$ serve the root.

\medskip
\noindent\textbf{Status.} \texttt{proved}; \texttt{af: validated} (T0). Workspace
\texttt{proofs/lem-hcb4-canonical-gram}: 9 nodes, all \texttt{validated}, root
\texttt{validated}, taint clean.

\noindent\textbf{Role in the chain.} Registry imports:
\texttt{lem-compcb-rectangular-product},
\texttt{lem-compcb-compressed-unit-norm},
\texttt{lem-compcb-corner-algebra} and
\texttt{lem-hcb3-uniform-square-lower}; the exported tree visibly consumes the first two,
the remaining two being available but unused in the validated derivation. This is the base
of the canonical layer of the H-CB tier. Canonical closeness
(\S\ref{sec:hcb4-canonical-closeness}) uses it once, to convert
$\lVert J_{P,Q,n}(Z)\rVert$ into $\nrm{Z}$; the canonical inverse estimate
(\S\ref{sec:hcb4-canonical-inverse}) uses its lower half twice, once for the bound
$\lVert J_n^{-1}\rVert\le\tfrac43$ and once for the norm equivalence that makes the
Neumann target complete. Through those two it reaches the canonical-identity clause of the
parent \texttt{conj-hcb} (\S\ref{sec:hcb}), which lists it among its sixteen dependencies.
Nothing here bears on \texttt{op-classical}, which is open.
